\section{Kết quả thực nghiệm}\label{sec:results}

\subsection{Môi trường và phiên bản}
\begin{table}[H]
\centering
\caption{Phiên bản phần mềm được ghi trong repo}
\label{tab:versions-vi}
\begin{tabular}{@{}ll@{}}
\toprule
\textbf{Thành phần} & \textbf{Phiên bản}\\
\midrule
K3s & \texttt{v1.35.5+k3s1}\\
Helm & \texttt{v4.1.4}\\
Argo CD & \texttt{v3.4.2}\\
Tailscale & \texttt{v1.98.2}\\
Redpanda & \texttt{v26.1.8}\\
RisingWave & \texttt{v2.8.3}\\
Vector & \texttt{0.55.0-alpine}\\
VictoriaMetrics & \texttt{v1.143.0}\\
Grafana & \texttt{v13.0.1+security-01}\\
cloudflared & \texttt{2026.5.0}\\
\bottomrule
\end{tabular}
\end{table}

\subsection{Cluster, GitOps và mức sẵn sàng pipeline}
\begin{table}[H]
\centering
\caption{Kết quả hạ tầng và pipeline}
\label{tab:infra-results-vi}
\begin{tabularx}{0.92\textwidth}{@{}p{0.48\textwidth}Y@{}}
\toprule
\textbf{Chỉ số} & \textbf{Kết quả}\\
\midrule
K3s nodes Ready & \texttt{3/3}\\
PVC Bound & \texttt{5/5}\\
Workloads Available sau replay & \texttt{23/23}\\
Argo CD applications chính & \texttt{Synced/Healthy}\\
RisingWave roles & meta, compute, compactor, frontend \texttt{RUNNING}\\
TLC lookup table & \texttt{265} dòng\\
Iceberg objects & data Parquet, equality-delete, position-delete, metadata\\
\bottomrule
\end{tabularx}
\end{table}

\subsection{Phân tích chất lượng dữ liệu nguồn}\label{sec:dataquality}
Trước khi diễn giải kết quả replay, đề tài đối chiếu chính xác bộ dữ liệu thô để giải thích vì sao logic green được thiết kế như đã trình bày. Toàn bộ tệp \path{data/raw/yellow_tripdata_2026-03.jsonl} được quét tuần tự một lần trên host \texttt{imac} bằng script Python, không lấy mẫu, để đếm các trường hợp vi phạm chất lượng và tính phân phối cơ bản của \texttt{fare\_amount}, \texttt{trip\_distance}, \texttt{pu\_location\_id} và \texttt{pickup\_time}. Kết quả được tổng hợp ở Bảng~\ref{tab:dataset-quality-vi}.

\begin{table}[H]
\centering
\caption{Phân tích chất lượng bộ dữ liệu NYC TLC Yellow Taxi tháng \texttt{2026-03} sau chuyển đổi JSONL}
\label{tab:dataset-quality-vi}
\begin{tabularx}{\textwidth}{@{}p{0.42\textwidth}Y@{}}
\toprule
\textbf{Chỉ số quan sát được trên bộ dữ liệu thô} & \textbf{Giá trị}\\
\midrule
Tổng số bản ghi JSONL & \texttt{3.952.451}\\
Bản ghi có \texttt{fare\_amount < 0} & \texttt{20.724} (\texttt{0,5243\%})\\
Bản ghi có \texttt{trip\_distance < 0} & \texttt{0}\\
Bản ghi có \texttt{fare\_amount = 0} & \texttt{2.995}\\
\texttt{fare\_amount} min / max & \texttt{-990,00} / \texttt{1.843,30} USD\\
\texttt{fare\_amount} mean / median & \texttt{21,33} / \texttt{15,60} USD\\
\texttt{trip\_distance} min / max & \texttt{0} / \texttt{288.381,68} miles\\
\texttt{trip\_distance} mean / median & \texttt{5,98} / \texttt{1,82} miles\\
\texttt{pu\_location\_id} xuất hiện trong bộ dữ liệu & \texttt{261}/\texttt{265} mã\\
Khoảng \texttt{pickup\_time} & \texttt{2008-12-31} đến \texttt{2026-04-01} (\texttt{35} ngày phân biệt)\\
Giờ cao điểm theo \texttt{pickup\_time} & \texttt{18:00} với \texttt{279.487} chuyến\\
Giờ thấp điểm theo \texttt{pickup\_time} & \texttt{04:00} với \texttt{35.646} chuyến\\
\bottomrule
\end{tabularx}
\end{table}

Kết quả này có ba ý nghĩa cho thiết kế và đánh giá:

\begin{itemize}
\item Bộ lọc \texttt{fare\_amount >= 0} của logic green được đặt trên một hiện tượng có thật: \texttt{20.724} bản ghi với cước âm, chiếm \texttt{0,5243\%} bộ dữ liệu thô. Tỉ lệ nhỏ này phù hợp với chênh lệch \texttt{8} chuyến giữa blue (\texttt{986}) và green (\texttt{978}) đo được trên mẫu nhỏ của lượt replay; ở mức \texttt{986} chuyến, kỳ vọng số bản ghi cước âm là khoảng $986 \times 0{,}5243\% \approx 5$ chuyến.
\item Không có bản ghi nào với \texttt{trip\_distance < 0} trong bộ dữ liệu; vì vậy predicate \texttt{trip\_distance >= 0} trong logic green là biện pháp phòng ngừa cho các tháng dữ liệu khác hoặc các bộ dữ liệu tương tự.
\item \texttt{trip\_distance} có giá trị tối đa \texttt{288.381,68} miles và \texttt{pickup\_time} chứa các mốc ngoài tháng \texttt{2026-03}, từ \texttt{2008-12-31} đến \texttt{2026-04-01}. Đây là các bất thường ở dữ liệu/schema thô chưa được logic green hiện tại xử lý và được liệt kê tại Mục~\ref{sec:limitations}.
\end{itemize}

Phân tích này cũng giải thích tại sao MV chỉ chứa \texttt{69} cặp (\texttt{borough}, \texttt{zone}) chứ không phải toàn bộ \texttt{261} mã \texttt{pu\_location\_id} mà bộ dữ liệu chạm vào: trong cửa sổ replay đã đo, Vector chỉ đưa được một mẫu \texttt{986} chuyến qua broker trước khi được dừng có chủ đích, và mẫu nhỏ này không phủ hết tất cả các vùng. Bối cảnh giờ cao điểm \texttt{18:00} và thấp điểm \texttt{04:00} minh họa tính chất tải biến thiên theo thời gian của ITS và củng cố luận điểm cần observability gần thời gian thực; tuy nhiên các con số này thuộc cấp bộ dữ liệu, không phải kết quả đo throughput của artifact.

\subsection{Replay ingest và đầu ra lakehouse}
Vector được mở trong khoảng có kiểm soát rồi đưa về \texttt{replicas=0}. File \path{05-mv-before-stop.txt} ghi public MV ở \texttt{66 zones / 912 trips}; sau khi các event đã nhận được xử lý hết, \path{05-mv-final-count.txt} ghi trạng thái ổn định \texttt{69 zones / 986 trips}. File \path{05-minio-iceberg-after-replay.txt} liệt kê đầu ra Iceberg mới, cho thấy đường ghi lakehouse đã hoạt động trong lượt replay.

Kết quả này không đại diện throughput tối đa hay xử lý toàn bộ \texttt{3.952.451} dòng. Các file cấu hình mức tải trong repo không kèm kết quả thực nghiệm, nên không được trình bày thành benchmark.

\begin{figure}[H]
\centering
\includegraphics[width=0.9\textwidth]{streaming-perf.png}
\caption{Dashboard hiệu năng streaming sau lượt replay/cutover. Consumer lag hiển thị \texttt{0} trong cửa sổ ảnh.}
\label{fig:streaming-vi}
\end{figure}

\subsection{Blue/Green MV cutover}
\begin{table}[H]
\centering
\caption{Kết quả Blue/Green MV swap}
\label{tab:cutover-results-vi}
\begin{tabularx}{0.92\textwidth}{@{}p{0.52\textwidth}Y@{}}
\toprule
\textbf{Chỉ số} & \textbf{Giá trị}\\
\midrule
Public MV trước swap & \texttt{69 zones / 986 trips}\\
Blue MV trước swap & \texttt{69 zones / 986 trips}\\
Green MV trước swap & \texttt{69 zones / 978 trips}\\
Cutover command & \texttt{ALTER MATERIALIZED VIEW ... SWAP WITH ...}\\
Thời gian do runbook ghi & \texttt{0,145226 s}\\
Lỗi query trong loop & \texttt{0}\\
RisingWave restarts trong swap & \texttt{0}\\
Public MV sau swap & \texttt{69 zones / 978 trips}\\
View mang tên green sau swap & \texttt{69 zones / 986 trips}\\
\bottomrule
\end{tabularx}
\end{table}

Chênh lệch \texttt{8} trips giữa blue và green phù hợp điều kiện green loại record có \texttt{fare\_amount < 0} hoặc \texttt{trip\_distance < 0}. Bằng chứng không cho thấy chênh lệch này phát sinh từ lỗi chuyển giao.

Query loop trong \path{06-query-loop-during-cutover.txt} ghi các truy vấn thành công khi public MV chuyển từ trạng thái \texttt{69|986} sang \texttt{69|978}, trong khoảng từ \texttt{16:17:36 UTC} đến \texttt{16:18:16 UTC} ngày \texttt{2026-05-22}; log không có dòng \texttt{ERROR}. File duration ghi timestamp swap \texttt{16:18:11 UTC} và duration \texttt{0,145226 s}; hai file VictoriaMetrics thu sau đó xác nhận metric duration và \texttt{continux\_query\_errors\_total = 0} đã được lưu. Do file swap/duration trong bundle nguồn có mốc cập nhật sau log chuyển trạng thái, số duration được dùng như phép đo runbook đã ghi nhận gắn với transition quan sát, chưa phải latency được kiểm chứng lại độc lập. Bằng chứng hỗ trợ nhận định không ghi nhận gián đoạn truy vấn, chứ không chứng minh mọi client hoặc mọi lỗi mạng đều vắng mặt.

\begin{figure}[H]
\centering
\includegraphics[width=0.9\textwidth]{cutover.png}
\caption{Dashboard cutover: green ready, query errors bằng \texttt{0} và restart bằng \texttt{0}; panel duration làm tròn giá trị nhỏ hơn một giây.}
\label{fig:cutover-vi}
\end{figure}

\subsection{Tính toàn vẹn quan sát được và tài nguyên}
Dashboard tính toàn vẹn dữ liệu hiển thị public MV có \texttt{69} dòng và checksum mismatch \texttt{1} sau swap. Exporter so public với logic baseline; sau cutover public mang green logic, nên mismatch là kết quả mong đợi của thay đổi định nghĩa, không phải bằng chứng duplicate hoặc loss. Panel \enquote{Iceberg freshness} hiển thị \enquote{56 years}; do metric timestamp chưa hợp lệ trong cấu hình quan sát, giá trị này không được dùng làm kết quả.

\begin{figure}[H]
\centering
\includegraphics[width=0.9\textwidth]{data-integrity.png}
\caption{Dashboard tính toàn vẹn dữ liệu sau swap; mismatch được diễn giải theo thay đổi logic.}
\label{fig:data-integrity-vi}
\end{figure}

Dashboard tài nguyên ghi workload availability \texttt{15/15} ở thời điểm ảnh và PVC MinIO/Redpanda khoảng \texttt{9,8\%}. Kết quả cục bộ trong \path{05-k3s-overview-after-replay.txt} ghi tình trạng toàn cluster sau replay là \texttt{23/23} workload available, RAM \texttt{imac} khoảng \texttt{49\%} và disk root khoảng \texttt{11\%}.

\begin{figure}[H]
\centering
\includegraphics[width=0.9\textwidth]{resource-util.png}
\caption{Dashboard resource utilization trong cửa sổ quan sát thực nghiệm.}
\label{fig:resource-vi}
\end{figure}

\subsection{Tổng hợp đánh giá}
\begin{table}[H]
\centering
\caption{Tổng hợp kết luận theo bằng chứng}
\begin{tabularx}{\textwidth}{@{}p{0.26\textwidth}p{0.34\textwidth}Y@{}}
\toprule
\textbf{Nhóm} & \textbf{Bằng chứng ghi nhận} & \textbf{Kết luận}\\
\midrule
Query-layer cutover & Runbook duration \texttt{0,145226 s}, lỗi \texttt{0}, restart \texttt{0} & Zero-downtime quan sát được cho public MV\\
Đầu ra lakehouse & Iceberg objects sau replay & Sink đã ghi đầu ra trong replay\\
Logic difference & Green ít hơn blue \texttt{8} trips & Thay đổi bộ lọc được quan sát\\
Cluster stability & \texttt{3/3}, \texttt{5/5}, \texttt{23/23} & Đủ chạy kịch bản đã đo\\
Exactly-Once end-to-end & Không có reconciliation & Không kết luận\\
Multi-load performance & Không có run results & Không kết luận\\
\bottomrule
\end{tabularx}
\end{table}
